\documentclass[11pt]{letter}
\usepackage[utf8]{inputenc}
\usepackage[T1]{fontenc}
\usepackage[margin=1in]{geometry}
\usepackage{amsmath}
\usepackage[colorlinks=true,linkcolor=red!70!black,citecolor=green!50!black,urlcolor=blue!80!black]{hyperref}

\signature{Mike Edison Almeida Vallejo\\\texttt{mike.almeida1721@gmail.com}}
\address{Mike Edison Almeida Vallejo \\ Independent Researcher}

\sloppy
\begin{document}

\begin{letter}{The Editorial Board\\Journal of Cosmology and Astroparticle Physics (JCAP)}

\opening{Dear Editors,}

I am pleased to submit the manuscript \textit{``Structural Self-Energy
Expansion: A Minimal-Parameter Dark-Energy Framework Anchored in
$\varphi$ and $\pi$''} for consideration in JCAP.

The paper addresses a central question in late-time cosmology: how much of the
dark-energy phenomenology can be fixed by structure rather than by free
parameters. We show that the \emph{shape} of cosmic acceleration---the
equation-of-state pair $(w_0,w_a)$, the canonical effective-field-theory
coefficients, and the inflationary spectral observables---follows algebraically
from two dimensionless constants, the golden ratio $\varphi$ and $\pi$, with no
tuned amplitude. In CMB-likelihood terms the framework is a $k=2$ fit: it fixes
four of $\Lambda$CDM's six parameters algebraically ($\Omega_b$, $\Omega_c$,
$n_s=1-\varphi^{-7}$, and the derived anchor $H_0/(\mathrm{km\,s^{-1}\,Mpc^{-1}})=3(\varphi+\pi)^2$), leaving only
the primordial amplitude $A_s$ and the optical depth $\tau$ ($k=2$ versus $k=6$).

The principal results that I believe merit publication are:

\begin{itemize}
  \item The bare ratios $w_0=-T_r/M_v=-0.839950$ and $w_a=-P_{\rm sc}/I_g=-0.669975$
  match the DESI DR2 $w_0w_a$CDM posterior at the $0.24\sigma$ level (Pantheon+; $0.2$--$1.8\sigma$ across SN compilations), with no
  fitted normalization.
  \item A Markov-chain analysis combining DESI BAO with the $\omega_m$-direct
  CMB anchor as prior yields $H_0=67.79\pm0.35~\mathrm{km\,s^{-1}\,Mpc^{-1}}$
  ($0.50\sigma$ from the algebraic anchor), and the
  full Planck PR4 \texttt{plik} likelihood (TT, TE, EE $+$ low-$\ell$ $+$ lensing)
  gives $\Delta\mathrm{BIC}=-32.9$ relative to $\Lambda$CDM ($k=2$ versus $k=6$),
  with a best-fit $\chi^2$ statistically indistinguishable from $\Lambda$CDM---the
  preference is driven by parsimony, not by a superior fit.
  \item The CMB matter density carries no fitted parameter: it follows from the
  forward identity $\omega_c=\mathrm{KAL}_0\,\omega_b\,n_s$, giving
  $\Omega_{m,\mathrm{CMB}}=\omega_m/h^2=0.308881$ in agreement with Planck at
  $0.88\sigma$ and reproducing the acoustic-peak structure. (The $0.41\sigma$
  quoted in an earlier version is the agreement of $\omega_c$ itself, not of
  $\Omega_{m,\mathrm{CMB}}$.)
  \item The inflationary sector predicts $n_s=1-\varphi^{-7}=0.965558$ and a
  tensor-to-scalar ratio $r=\varphi^{-10}=8.130619\times10^{-3}$, falsifiable by
  LiteBIRD and CMB-S4.
\end{itemize}

In keeping with the standards of the journal, the manuscript is explicit about
the framework's limits. The absolute vacuum scale enters as a single measured
input rather than a derived quantity; the non-linear closure of the growth
amplitude is deferred to N-body work; and the mechanism behind MIRA remains
an open problem rather than a claimed result. \emph{(Updated 2026-09-07:
earlier versions of this letter said $S_8$ ``remains in tension with
weak-lensing surveys, relaxed only in the conditional two-sector extension'',
and cited the $\varphi$-dark-matter mass coefficient. Both are withdrawn.
Against the raw KiDS-1000 data with $A_s$ free --- it is a free parameter of
the model --- the converged MCMC gives $S_8=0.7555\pm0.0192$, which is
$0.11\sigma$ from KiDS: there is no tension to relax. The $3.5\sigma$ figure
was an artefact of fixing $A_s$ to Planck, i.e.\ of importing the
Planck--KiDS tension. The two-sector extension and the particle were retired
on 2026-08-01.)} I have deliberately separated the algebraically fixed core from these
ongoing questions so that referees can assess each on its own terms.

The work is original, has not been published elsewhere, and is not under
consideration by any other journal. I have no competing interests to declare.
All numerical results are reproducible from the analysis scripts accompanying
the manuscript.

Thank you for your time and consideration. I would be glad to respond to any
questions from the editors or referees.

\closing{Sincerely,}

\end{letter}
\end{document}
